% AI-written, AI cross-reviewed; no human sign-off; not a claim that RH is proved
% Atlas unit P027. Local manuscript v3, item 7.4, heading on printed PDF p. 74.
% UpperTheory.V.Alg.M.main_holds
% paper2/inputs_one.tex:131–203. Verbatim excerpt; SHA-256 4022491f2705cd8ce02fe1d6987fd69494872a2e1f32b30a99fe03e40b520cff.
% Original TeX uses the macros and cross-reference keys of the source manuscript.

\begin{formalstatement}[Length measured by two-sided orbits]\label{pn:P027}
Let \(O\) be a complete characteristic-zero discrete valuation ring and
\(K\) an algebraically closed extension of its fraction field with
valuation \(v\) extending the normalized order. Let \(A=\bigoplus_j A_j\)
be graded, with marked \(O\)-point \(x\) zero on nonzero weights.
Its degree-zero algebra is Noetherian, complete local, with local
structure map surjective on residue fields. It generates \(A\) together
with finitely many elements of weights \(-2,-1,1,2\).
Let \(I\) be generated by lists of weights \((-1,-2,1,2,0)\) and sizes
\((r+1,s,r',s',m)\); denote the first list by \(F\), and assume the
weight-zero equations vanish at \(x\). Put
\[
 P=\{b\in A:\,tb\in(\ker x)^2\text{ for some }0\ne t\in O\},
 \qquad L_j=\operatorname{im}(A_j\cap\ker x\longrightarrow A/P).
\]
Assume \(r\ge1\) and the smooth graded chart \(M_1\): its variable
counts in the five weights are \((r,s,r'+1,s'+1,m)\); individual
linear coordinates can be isolated, and vanishing of the linear
part on \(\ker x\) implies membership in \(P\). The negative and
weight-zero equation lists span the corresponding cotangent lattices
over the fraction field. Assume the image of \(A_0\) in \(A/I\)
is finite over \(O\). In what follows take \(K\)-points of \(A/I\).
For each point vanishing on both signs other
than \(x_K\), at least one strictly one-sided sign is absent in its
fibre over \(A_0\). At every two-sided point (nonzero on some element
of each sign), point derivations killing \(I\) have dimension at most
one. Let \(a\in A_0\) vanish at \(x\) and be nonzero at all two-sided
points. There exists a finite set of two-sided points meeting every
such \(K^\times\)-orbit. For every such cover \(T\),
\[
 L=\operatorname{length}_O(L_{-1}/\langle\overline F_i\rangle)
 <\infty,\qquad L\le 2\sum_{\xi\in T}v(\xi(a)).
\]
\end{formalstatement}\nodestatus{P027}
\begin{outline}
We recall the two length estimates in Blueprint Z1.4~\bploc{sec:P-z14-full}, \emph{Graded
cotangent length}. After saturation the cotangent lattice in weight
\(-1\) is free of rank \(r\). Smith reduction of the first list gives
elementary divisors \(\pi^{c_i}\), hence \(L=\sum_i c_i\); the changed
equations generate the same ideal. Transverse slices to the
one-dimensional torus orbits are zero-dimensional by the derivation
bound. Together with finiteness in degree zero this yields a finite cover.
In the multigraded Rees calculation let \(\chi_d\) denote the eventual
Euler characteristic of the cut sections modulo \(a\) in multidegree \(d\).

On the chart the leading negative equations have the form
\(\pi^{2c_i}y_i^2\). Removing their powers of \(\pi\) gives
\(2\sum_i c_i\) times the mixed difference of the reduced Hilbert
function. Its eventual value is \(2^{r-1+r'}\): the generating
function is
\[
 (1+X)^{r-1}(1+Y)^{r'}\frac{1-X^4Y^2}{(1-X)(1-Y)(1-Y^2)}.
\]
For high \(X\)-degrees both shifts contribute one, cancelling the
factor \(1-Y^2\). Nonnegativity of the remaining Euler characteristic
gives \(2^{r+r'}L\le\chi_d\).

Placing \(a\) last bounds \(\chi_d\) by the Euler characteristic of
multiplication by \(a\) on the successive quotient. Its finite
degree-\(d\) module admits a prime filtration: the vertical terms
cancel and a horizontal term contributes its generic length times
\(v(\xi(a))\). The one-sided condition places the contributing
embeddings among the restrictions of \(T\). Before squaring the
equations, each localized weight piece has length at most one.
Indeed, the transverse slice is reduced, and multiplying by an
invertible pair of opposite weights makes every nonzero piece cyclic
over its residue field. Replacing a section by its square at most
doubles the bound, by filtering the squared quotient by the image of
the section. There are \(r+1+r'\) such replacements. Thus
\(\chi_d\le 2^{r+1+r'}\sum_{\xi\in T}v(\xi(a))\).
The values are nonnegative, so repeated restrictions cause no problem.
Cancelling \(2^{r+r'}\) also covers an empty \(T\).
\end{outline}
